\scanpage{35}
Cela ne suffit pas pour bien des questions : on aura besoin de partitions privilégiées qui \deleted{\unsure{suff}} soient plus fines qu'une partition \inserted{admissible} donnée, \inserted{\deleted{\unsure{ou}}(\unsure{plus gross.})} ou telles que \deleted{\unsure{donnés}} pour des sous-espaces fermés \uline{donnés} \inserted{$X_i$} de $X$ en nb fini, chacun soit réunion de feuilles et que la filtration induite soit \deleted{équisingulière}\inserted{privilégiée}. Ou pour une application donnée $X\to Y$, on voudrait trouver simultanément des partitions privilégiées compatibles avec cette application. Tous ces cas \unsure{sembleraient|semblent avoir} couverts par l'énoncé conjectural suivant :

\uline{Conjecture 2} Soient $I$ un ens. ordonné \uuline{fini}, $(X_i)_{i\in I}$ un diagramme de type $I$ d'espaces, avec les applications $X_i\to X_j$ \uline{propres}. Alors il existe des partitions privilégiées de ce diagramme, et il y a parmi elles une moins fine (appelée comme de juste ``canonique'').

\uline{Remarque} Soit une \deleted{\unsure{filtr}} partition admissible pour un diagramme $(X_i)_{i\in I}$. Si $X_i\overset{f}{\unsure{-|\to}}X_j$ dans ce diagramme est \uline{propre}, on voit aussitôt que pour tte feuille $F_{\deleted{\unsure{\alpha}}}$ de $X_i$, s'appliquant \inserted{donc} dans une feuille $\Phi$ de $X_j$, l'image est $\Phi$ tout entier (et non seulement un ouvert de $\Phi$) \deleted{\unsure{car l'image}}\inserted{car l'image de} \deleted{\illegible{3}} $\overline F\cap f^{-1}(\Phi)$ est à la fois ouvert \inserted{dans $\Phi$}\deleted{\unsure{plus}}
\mspage{36}
